\documentclass[11pt]{article}
\usepackage[margin=1in]{geometry}
\usepackage{booktabs}
\usepackage{amsmath}
\usepackage{amssymb}
\usepackage{graphicx}
\usepackage[colorlinks=true,linkcolor=blue,citecolor=blue]{hyperref}

\renewcommand{\thetable}{S\arabic{table}}
\renewcommand{\thesection}{S\arabic{section}}

\title{Supplementary Material\\
\large Deep Neural Network Surrogates for Approximate Bayesian Computation:\\
A Simulation Study for Mutation-Rate Estimation\\
in Markov Branching Processes}
\author{}
\date{}

\begin{document}
\maketitle

\noindent \textbf{Journal:} Statistics and Computing (submitted).
\textbf{Corresponding author:} [PLACEHOLDER: name, affiliation and email,
to match the main manuscript's title page once the author list is
finalised.]

\noindent This document contains the full architecture-, activation-, and
capacity-search tables referenced from the main text as ``Supplementary
Material S1'', together with the full per-family results of Study II's
architecture comparison. Each table below corresponds to a numbered claim in
the main text; the surrounding prose there states the headline conclusion,
and is not repeated here except as needed to make each table self-contained.

\section{Study I: the full activation, normalisation and capacity search}
\label{ssec:s1-1d-arch}

The deployed one-dimensional surrogate (Section~5.3 of the main text) is a
shallow funnel multilayer perceptron, input $\theta=\log_{10}p$, two
GELU-activated dense layers at width $32$--$16$ ($626$ parameters). That
choice is the outcome of the two-stage search reported in full here:
activation/normalisation first, at a $128$--$64$ baseline, then capacity,
once the activation question was settled.

\begin{table}[h]
\centering
\caption{Architecture search, one-dimensional model: mean-curve fit MSE
against the denoised response curve (irreducible replicate-noise floor
$\approx 4.7\times 10^{-3}$).}
\label{stab:arch1D}
\resizebox{\linewidth}{!}{%
\begin{tabular}{lrr}
\toprule
Model & Mean-curve MSE & vs.\ GP \\
\midrule
Gaussian process (GPS-ABC baseline) & $3.836\times10^{-4}$ & --- \\
\textbf{DNN, GELU (128, 64), no BatchNorm [selects activation]} & $\mathbf{3.888\times10^{-4}}$ & $+1\%$ \\
DNN ensemble $\times 5$, SiLU (128, 128, 64) & $3.960\times10^{-4}$ & $+3\%$ \\
DNN ensemble $\times 5$, tanh (256, 128) & $4.178\times10^{-4}$ & $+9\%$ \\
DNN, SiLU (128, 128, 64) & $4.400\times10^{-4}$ & $+15\%$ \\
DNN, tanh (64, 64, 32) & $4.435\times10^{-4}$ & $+16\%$ \\
DNN, tanh (256, 128) & $4.503\times10^{-4}$ & $+17\%$ \\
DNN, ReLU + BatchNorm + dropout (64, 64, 32) & $4.352\times10^{-3}$ & $+1035\%$ \\
\bottomrule
\end{tabular}%
}
\end{table}

The row that determines every downstream comparison in the main text is the
last one: an otherwise-reasonable network carrying BatchNorm and dropout fits
\emph{eleven times worse} than the same network without them. BatchNorm
normalises per-minibatch statistics, useful in deep classification networks;
but on a smooth scalar regression over a one-sided domain (the $\theta$-grid
has a hard edge at $-8$) it injects batch-dependent noise and biases
predictions near that boundary. A follow-up ablation isolating only the
activation function (same 128--64 architecture, no BatchNorm, averaged over
three random seeds) found ReLU, SiLU, GELU, and tanh within $7\%$ of one
another ($3.95$--$4.23\times10^{-4}$): a statistical tie, with ReLU
marginally ahead. Because accuracy does not separate them, the choice of
GELU is made on structural grounds (smoothness, immunity to the dead-neuron
failure mode) at no measurable accuracy cost. A second search round using
$7$-network deep ensembles (mean-curve MSE $3.94$--$4.25\times10^{-4}$
across four ensemble configurations) did not improve on the single chosen
network, so the simpler model is retained throughout.

\begin{table}[h]
\centering
\caption{One-dimensional model: does capacity matter? Curve MSE is
against the denoised response curve, mean over three seeds. Downstream
MSE and 95\% CI length are DNN-ABC at the same scale as the main text's
mean-squared-error table ($40$ replicates/cell, full $p\times J$ grid), one
seed per row.}
\label{stab:capacity1D}
\setlength{\tabcolsep}{4.5pt}
\resizebox{\linewidth}{!}{%
\begin{tabular}{lrrrr}
\toprule
Hidden layers & Parameters & Curve MSE ($\times$ GP) & Downstream MSE & Downstream 95\% CI \\
\midrule
GP (GPS-ABC baseline) & --- & $3.836\times10^{-4}$ (---) & $2.702\times10^{-6}$ & $2.137\times10^{-3}$ \\
\textbf{32--16 (deployed)} & \textbf{626} & $\mathbf{3.896\times10^{-4}}$ ($\mathbf{+1.6\%}$) & $\mathbf{1.433\times10^{-6}}$ & $\mathbf{1.300\times10^{-3}}$ \\
128--64 (original) & $8{,}642$ & $4.129\times10^{-4}$ ($+7.6\%$) & $1.514\times10^{-6}$ & $1.308\times10^{-3}$ \\
64--32 & $2{,}274$ & $3.957\times10^{-4}$ ($+3.2\%$) & $1.565\times10^{-6}$ & $1.377\times10^{-3}$ \\
16 (1 layer) & $66$ & $4.460\times10^{-4}$ ($+16.3\%$) & $1.579\times10^{-6}$ & $1.375\times10^{-3}$ \\
8 (1 layer) & $34$ & $4.983\times10^{-4}$ ($+29.9\%$) & $1.668\times10^{-6}$ & $1.393\times10^{-3}$ \\
16--8 & $186$ & $4.508\times10^{-4}$ ($+17.5\%$) & $1.810\times10^{-6}$ & $1.528\times10^{-3}$ \\
Linear (control) & $4$ & $4.672\times10^{-4}$ ($+21.8\%$) & --- & --- \\
\bottomrule
\end{tabular}%
}
\end{table}

On raw curve fit, capacity is not merely unnecessary above a small
threshold: it is mildly \emph{counter-productive} at the original size:
$128$--$64$ has the largest gap to the GP of any two-layer candidate tested
($+7.6\%$), while $32$--$16$, at $14\times$ fewer parameters, comes closest to
matching it ($+1.6\%$). On the downstream task the surrogate is actually
used for, $32$--$16$ is also the best point estimate on both accuracy and
interval length of every candidate tested, $128$--$64$ included, though
attaching Monte Carlo standard errors to the gap shows this specific edge is
not itself resolved (all nine $\Delta/\mathrm{SE}$ values fall below $0.6$ in
magnitude). Nothing in the data argues against the smaller network either:
$32$--$16$ never does measurably worse than $128$--$64$ in any cell. Only the
linear control shows a clear, unambiguous cost, confirming some nonlinearity
is required, just not much capacity beyond that.

\section{Study II: the full activation search}
\label{ssec:s1-3d-act}

The deployed three-dimensional surrogate uses GELU into tanh (different
activations per hidden layer) at width $64\to32$. Ten activations spanning
three families (piecewise-linear: ReLU, LeakyReLU, PReLU; smooth
saturating: tanh, ELU, SELU; and smooth unbounded: softplus, GELU, SiLU,
Mish) were compared on the deployed network, five seeds each.

\begin{table}[h]
\centering
\caption{Study II: activation comparison on the deployed $64\to32$ network,
five seeds each. MSE is against held-out design-point means of
$\log_{10}S$; the irreducible floor is $3.451\times10^{-4}$. The final
column compares each row to the best, in units of the standard error of the
difference. Standard deviations are across seeds (sample, $n-1$).}
\label{stab:act3D}
\begin{tabular}{llrrr}
\toprule
Activation & Family & $\times$ floor & $\pm$ sd (seeds) & $t$ vs.\ best \\
\midrule
\textbf{GELU} & smooth unbounded & \textbf{1.072} & $9.5\times10^{-6}$ & --- \\
tanh & smooth saturating & 1.078 & $6.0\times10^{-6}$ & 0.42 \\
LeakyReLU & piecewise-linear & 1.080 & $6.4\times10^{-6}$ & 0.50 \\
ReLU & piecewise-linear & 1.082 & $9.5\times10^{-6}$ & 0.54 \\
Mish & smooth unbounded & 1.094 & $7.2\times10^{-6}$ & 1.38 \\
SiLU & smooth unbounded & 1.130 & $1.2\times10^{-5}$ & 3.00 \\
PReLU & piecewise-linear & 1.135 & $1.3\times10^{-5}$ & 3.04 \\
ELU & smooth saturating & 1.166 & $2.3\times10^{-5}$ & 2.95 \\
SELU & smooth saturating & 1.265 & $3.3\times10^{-5}$ & 4.35 \\
softplus & smooth unbounded & 1.383 & $8.7\times10^{-5}$ & 2.74 \\
\bottomrule
\end{tabular}
\end{table}

Two results follow, pulling in opposite directions. First, the choice
\emph{can} matter: the best-to-worst spread is $29\%$ of the best row's MSE
at $t=2.7$, outside seed noise, and softplus fails outright (strictly
positive, so every downstream layer must undo an accumulated offset).
Second, the five leading candidates (GELU, tanh, LeakyReLU, ReLU and
Mish) are mutually indistinguishable, every pairwise $|t|$ among them
below $2$. SELU's poor placement reflects that it presupposes LeCun-normal
initialisation and AlphaDropout, absent under this project's standard
initialisation; its position here is evidence about this configuration, not
a general verdict on SELU.

Because the two hidden layers need not share an activation, all
$10\times10=100$ ordered per-layer assignments were also compared. All $100$
pairs are screened; the ten leading candidates are then re-fit on fifteen
\emph{fresh} seeds that played no part in selection, with the test split
touched exactly once, for the eventual winner alone.

\begin{table}[h]
\centering
\caption{Study II: activation selection with no test-set leakage. All $100$
ordered pairs were screened on validation; the ten finalists were re-fit on
fifteen fresh validation seeds. Selection is on validation MSE; the test
column is shown for reference and played no part in the choice. Test
$\times$ floor is against the irreducible test floor $3.451\times10^{-4}$.}
\label{stab:actsel3D}
\begin{tabular}{llrrrr}
\toprule
Layer 1 (64) & Layer 2 (32) & Val MSE & $\pm$ sd & $t$ vs.\ best & Test $\times$ floor \\
\midrule
\textbf{GELU} & \textbf{tanh} & $\mathbf{2.980\times10^{-4}}$ & $1.3\times10^{-5}$ & --- & 1.062 \\
Mish & tanh & $3.006\times10^{-4}$ & $1.4\times10^{-5}$ & 0.54 & 1.076 \\
SiLU & tanh & $3.021\times10^{-4}$ & $1.4\times10^{-5}$ & 0.86 & 1.086 \\
LeakyReLU & LeakyReLU & $3.059\times10^{-4}$ & $8.4\times10^{-6}$ & 2.00 & 1.075 \\
PReLU & ReLU & $3.071\times10^{-4}$ & $1.2\times10^{-5}$ & 2.05 & 1.077 \\
Mish & Mish & $3.120\times10^{-4}$ & $1.1\times10^{-5}$ & 3.19 & 1.080 \\
ELU & GELU & $3.137\times10^{-4}$ & $8.5\times10^{-6}$ & 3.99 & 1.093 \\
SiLU & Mish & $3.164\times10^{-4}$ & $1.4\times10^{-5}$ & 3.84 & 1.097 \\
GELU & Mish & $3.172\times10^{-4}$ & $1.7\times10^{-5}$ & 3.51 & 1.086 \\
Mish & SiLU & $3.205\times10^{-4}$ & $1.8\times10^{-5}$ & 3.94 & 1.101 \\
\bottomrule
\end{tabular}
\end{table}

The selection optimism this two-stage discipline guards against is not
small: the finalists perform, on average, $1.03\times10^{-5}$ \emph{worse}
on the fresh seeds than on the seeds that chose them, $45\%$ of the entire
spread of the confirmation table and comparable to the across-seed standard
deviation itself. The pair that screened best, SiLU into tanh, finished
third on confirmation: a single-stage search would have reported that
movement as a real result. Under the fourth-root target used throughout the
main text, GELU into tanh is best on both validation and test, so scoring
selection on test would not, here, have changed the choice; under the
square-root target tried during development it did, which is why the
two-stage, validation-only discipline is treated as a requirement rather
than a formality.

\section{Study II: the architecture-family comparison}
\label{ssec:s1-3d-family}

A 1-D CNN (two convolutional layers, kernel width 3), a GRU, and an LSTM
(each unrolling the three inputs $(\log_{10}p_1,\log_{10}p_2,\tau)$ as three
timesteps), with comparable parameter counts ($2{,}402$--$3{,}426$) to the
deployed FFN, were each evaluated exactly as GPS-ABC and DNN-ABC were: a
surrogate-fit comparison against the irreducible noise floor, and the
identical full ABC-MCMC recovery task (same three truth triples, $16$
replicates, $3{,}000$ iterations, $1{,}000$ burn-in).

\begin{table}[h]
\centering
\caption{Surrogate fit quality by architecture family, test set. The
irreducible noise floor is $3.451\times10^{-4}$. The deployed FFN's row is
read from the capacity study rather than restated, so it necessarily
describes the same target as the rows beside it.}
\label{stab:archfam3D}
\resizebox{\linewidth}{!}{%
\begin{tabular}{lrrrrr}
\toprule
Architecture & Parameters & MSE (design means) & $\times$ floor & $R^2$ & 95\% coverage \\
\midrule
CNN1D (2 conv layers, $k{=}3$) & 2{,}482 & $3.589\times10^{-4}$ & 1.04 & 0.99670 & 0.949 \\
\textbf{FFN 64--32 [deployed]} & \textbf{2{,}402} & $\mathbf{3.749\times10^{-4}}$ & \textbf{1.09} & \textbf{0.99656} & \textbf{0.952} \\
LSTM (hidden $24$) & 2{,}642 & $4.221\times10^{-4}$ & 1.22 & 0.99612 & 0.950 \\
RNN (GRU, hidden $32$) & 3{,}426 & $4.520\times10^{-4}$ & 1.31 & 0.99585 & 0.949 \\
\bottomrule
\end{tabular}%
}
\end{table}

\begin{table}[h]
\centering
\caption{ABC-MCMC recovery of $(p_1,p_2,\tau)$ by architecture family,
per-parameter RMSE (log$_{10}$ scale for $p_1,p_2$; absolute time units for
$\tau$), averaged over the three truth triples used throughout Study II.
Bold marks the best value in each row.}
\label{stab:archfam3D_abc}
\begin{tabular}{lrrrrr}
\toprule
Parameter & GPS-ABC & DNN-ABC (FFN) & CNN1D-ABC & RNN-ABC & LSTM-ABC \\
\midrule
$p_1$ RMSE & 1.391 & 1.407 & \textbf{1.373} & 1.378 & 1.389 \\
$p_1$ coverage & 1.000 & 0.938 & 0.958 & 0.917 & 0.979 \\
$p_2$ RMSE & \textbf{0.156} & 0.258 & 0.203 & 0.312 & 0.209 \\
$p_2$ coverage & 1.000 & 0.979 & 0.979 & 0.938 & 0.958 \\
$\tau$ RMSE & \textbf{1.422} & 1.753 & 1.688 & 1.700 & 1.631 \\
$\tau$ coverage & 1.000 & 1.000 & 1.000 & 1.000 & 1.000 \\
\bottomrule
\end{tabular}
\end{table}

CNN1D statistically ties the deployed FFN ($4.3\%$ lower MSE, inside the
seed-to-seed spread already documented for the capacity study), but costs
$143$\,\textmu s per query against the FFN's $40$, the wrong trade for a tie.
LSTM and RNN are measurably worse ($+12.6\%$, $+20.6\%$ MSE; both gaps an
order of magnitude larger than their own across-seed standard deviation).
Attaching Monte Carlo standard errors to every one of the $3\times3\times3=27$
new-family, truth and parameter cells, and comparing each new family against
both GPS-ABC and DNN-ABC at the $2$-MCSE level ($54$ pairwise comparisons in
total), gives a picture consistent with the fit-quality result, and a
slightly starker one: $49$ of $54$ ($91\%$) are statistical ties. All $5$
resolved comparisons favour GPS-ABC, and \emph{none} favours a new family
anywhere. Three of the five are CNN1D-ABC and two RNN-ABC; LSTM-ABC is tied
in every one of its eighteen comparisons.

Interval calibration adds nothing to the picture either way. Pooled coverage
is $0.951$--$0.979$ for the four neural families and $1.000$ for both
Gaussian processes, against a nominal $0.95$. An earlier version of this
comparison read a narrative into that column (the FFN under-covering,
CNN1D inheriting and worsening it, LSTM correcting it) which the corrected
numbers do not support: the spread across all four neural families is under
three percentage points, which at sixteen replicates is at most one cell per
parameter. Every method over-covers on $\tau$ at exactly $1.000$, reflecting
that parameter's weak identifiability rather than any property of an
architecture.

\end{document}
